\chapter{Cơ sở lý thuyết và tổng quan nghiên cứu}

\section{Dữ liệu log hệ thống}
Log hệ thống là chuỗi bản ghi mô tả trạng thái, hành vi hoặc sự kiện của phần
mềm và hạ tầng. Một bản ghi thường chứa timestamp, mức severity, component,
message và tham số thay đổi.

\subsection{Đặc điểm của log}
Log bán cấu trúc, phi dừng, có dung lượng lớn và mất cân bằng lớp. Template có
thể thay đổi khi phần mềm phát triển; một template chưa từng thấy không tự động
là bất thường. Nhãn có thể gắn với một dòng, block hoặc cửa sổ thời gian.

\subsection{Log parsing bằng Drain}
Drain chuyển message thô thành template ổn định bằng cây có độ sâu cố định,
ngưỡng tương đồng và cơ chế masking tham số. Template giúp gom các log cùng cấu
trúc thành event ID. Lỗi parser có thể lan truyền sang biểu diễn và mô hình
phát hiện.

\subsection{Gom nhóm log}
Log được gom theo block/session, fixed-size window, sliding window hoặc trace.
HDFS\_v1 gom theo block ID, trong khi BGL có thể gom theo timestamp hoặc đánh
giá ở cấp dòng tùy giao thức.

\section{Bộ dữ liệu nghiên cứu}
\subsection{LogHub}
LogHub là bộ dữ liệu công khai tổng hợp nhiều log hệ thống và cung cấp tài
nguyên cho nghiên cứu phân tích log.

\subsection{HDFS\_v1}
HDFS\_v1 chứa 11.175.629 dòng log thô. Nhãn được gắn ở cấp block thông qua
\texttt{anomaly\_label.csv}: 558.223 block Normal và 16.838 block Anomaly.
Tổng số block là 575.061.

\subsection{BGL}
BGL chứa 4.747.963 dòng log thô trong dữ liệu tải về. File raw hiện không
cung cấp ground-truth label theo từng dòng. Do đó nghiên cứu cần thêm nguồn
nhãn chuẩn trước khi tính chỉ số đánh giá BGL.
